\section{Dinámica de Newton--Euler completa}
\label{sec:dinamica}

\subsection{Estado y entrada}

El vector de estado tiene 17 componentes (\cod{dynamics.py}, índices
\cod{IP, IV, IQ, IW, IR}):
\begin{equation}
\vect x = \big[\ \vect p\T\ \ \vect v\T\ \ \vect q\T\ \ \vect\omega\T\ \ \Omega_1\ \ \Omega_2\ \ \Omega_3\ \ \Omega_4\ \big]\T\in\mathbb R^{3+3+4+3+4=17} .
\label{eq:estado}
\end{equation}
La entrada son los cuatro aceleradores $\sigma_i\in[0,1]$, retenidos
(orden cero) durante cada paso de integración. Las perturbaciones entran por
un objeto \cod{Environment}: velocidad del viento $\vect v_w$ en $\E$, fuerza
externa $\vect F_{ext}$ en $\E$, momento externo $\vect\tau_{ext}$ en $\B$,
eficiencia de cada motor $\eta_i$ y un indicador de efecto suelo.

Incluir las velocidades de los rotores en el estado, en lugar de suponer que
el empuje sigue instantáneamente al mando, es la diferencia entre un modelo
útil para sintonizar y uno que no lo es. El rezago del motor ($T_m=20$~ms en
el PAE-X450) es el mayor retardo de la cadena de control de actitud; sin él
cualquier ganancia de tasa parece estable.

\subsection{Supuestos}

\begin{supuesto}[Supuestos del modelo de sólido rígido]
\begin{itemize}[leftmargin=1.2em]
\item Tierra plana y no rotante: $\E$ es inercial y la gravedad es uniforme.
A las velocidades y distancias de un multirrotor de laboratorio, la
aceleración de Coriolis terrestre ($\sim10^{-4}$~m/s$^2$) es despreciable.
\item Cuerpo rígido con masa constante (salvo los eventos de carga útil,
\cref{sec:perturbaciones}) y tensor de inercia diagonal (productos de inercia
nulos por simetría).
\item Rotores en el plano horizontal que contiene al centro de masa: el
empuje no genera momento por un brazo vertical.
\item Las fuerzas aerodinámicas del fuselaje actúan en el centro de masa.
\end{itemize}
\end{supuesto}

\subsection{Traslación (Newton, en el marco inercial)}

\begin{equation}
m\,\dot{\vect v} = m g\,\eu + \R(\vect q)\,\vect f_B + \vect F_{ext},
\qquad
\vect f_B = \begin{bmatrix} f_{a,x}\\ f_{a,y}\\ f_{a,z} - \sum_{i=1}^4 T_i\end{bmatrix},
\label{eq:newton}
\end{equation}
con $T_i = C_T\Omega_i^2$ (multiplicado por el factor de efecto suelo cuando
está activo) y $\vect f_a$ la fuerza aerodinámica del fuselaje en $\B$
(\cref{sec:aero}). La ecuación se escribe en $\E$ porque ahí la gravedad es
constante y la posición se integra directamente; la fuerza de los rotores
se escribe en $\B$ porque ahí es constante en dirección. La matriz $\R$ hace
el puente, y con ella entra el acoplamiento que la 2.ª edición no tenía:
inclinar el dron inclina el empuje, y el dron se desplaza.

\subsection{Rotación (Euler, en el marco cuerpo)}

El momento angular total del vehículo es el del cuerpo más el de los cuatro
rotores. Si cada rotor gira con velocidad $\Omega_i$ alrededor de $z_B$ en el
sentido $e_i$, un rotor antihorario visto desde arriba tiene su vector de
velocidad angular hacia arriba, es decir, hacia $-z_B$. El momento angular
de los rotores es entonces
\begin{equation}
\vect h_r = \begin{bmatrix}0\\0\\h_z\end{bmatrix},\qquad
h_z = -J_m\sum_{i=1}^4 e_i\,\Omega_i .
\label{eq:hr}
\end{equation}
Derivando $J\vect\omega+\vect h_r$ en el marco móvil e igualando al momento
externo se obtiene
\begin{equation}
J\dot{\vect\omega} = \vect\tau_B - \vect\omega\times(J\vect\omega)
 - \vect\omega\times\vect h_r - \underset{\scriptscriptstyle \text{omitido}}{\dot{\vect h}_r}
 + \vect\tau_a + \vect\tau_{ext},
\label{eq:euler}
\end{equation}
donde el momento de los rotores es
\begin{equation}
\vect\tau_B=\begin{bmatrix}\tau_x\\ \tau_y\\ \tau_z\end{bmatrix}
=\begin{bmatrix}
-\sum_i y_i\,T_i\\
\phantom{-}\sum_i x_i\,T_i\\
\phantom{-}\sum_i e_i\,Q_i
\end{bmatrix},\qquad Q_i=C_M\Omega_i^2 .
\label{eq:tauB}
\end{equation}
El signo de la fila de roll sale del producto cruz: el rotor $i$ está en
$(x_i,y_i,0)$ y empuja con $(0,0,-T_i)$, así que
$(x_i,y_i,0)\times(0,0,-T_i)=(-y_iT_i,\ x_iT_i,\ 0)$. Un rotor a la derecha
($y_i>0$) que empuja más produce roll negativo (levanta el ala derecha).

En componentes, con $J$ diagonal, el código (\cod{QuadrotorDynamics.derivative})
evalúa
\begin{align}
J_{xx}\dot p &= \tau_x - (J_{zz}-J_{yy})\,q r - q\,h_z + \tau_{a,x} + \tau_{ext,x},\nonumber\\
J_{yy}\dot q &= \tau_y - (J_{xx}-J_{zz})\,r p + p\,h_z + \tau_{a,y} + \tau_{ext,y},\label{eq:euler_comp}\\
J_{zz}\dot r &= \tau_z - (J_{yy}-J_{xx})\,p q + \tau_{a,z} + \tau_{ext,z}.\nonumber
\end{align}

\paragraph{El término giroscópico de los rotores, $-\vect\omega\times\vect h_r$.}
Solo afecta a roll y pitch, y es proporcional a la \emph{diferencia} entre
las velocidades de los rotores CCW y CW. En hover esa diferencia es cero y el
término desaparece; aparece cuando hay mando de guiñada mientras el dron
gira en roll o pitch. Con $J_m/J_{zz}=\val{JmJz}$ su efecto es pequeño en el
PAE-X450, pero se incluye porque no cuesta nada y está en todos los modelos
de referencia \cite{qq17,hb11}.

\subsection{El término de reacción por aceleración de los rotores, omitido}
\label{sec:reaccion}

El término $-\dot{\vect h}_r = J_m\sum_i e_i\dot\Omega_i\,\vect e_3$ es la
reacción que siente el cuerpo cuando los motores \emph{aceleran} sus rotores.
Está en las ecuaciones de Phang et al.\ \cite{hb11} (ecs.~11.14--11.16), y el
código lo calcula, pero no lo aplica, salvo que la dinámica se construya
con la opción \cod{include\_rotor\_reaction=True}. La decisión no es
de comodidad; es de consistencia con el resto del modelo.

Con el propulsor de primer orden de la \cref{sec:propulsion}, un escalón del
acelerador desde hover hasta el máximo produce, en el primer instante,
\begin{equation}
\dot\Omega = \frac{\Omega_{max}-\Omega_h}{T_m}
= \frac{\val{wmax}-\val{wh}}{0.02} = \val{dOm}\ \text{rad/s}^2,
\qquad
J_m\dot\Omega = \val{reac}\ \text{N·m por motor}.
\end{equation}
Ese par tendría que darlo el motor. Un motor 2212 de 920~KV (el típico de un
F450) tiene $K_t = 60/(2\pi\cdot920) = \val{Kt}$~N·m/A; con un límite de
corriente de 15~A (supuesto razonable para un ESC de 20--30~A) entrega como
mucho $\val{taumot}$~N·m, unas \val{reacRatio} veces menos. El par
aerodinámico máximo de la hélice es $\val{Qmax}$~N·m, \val{reacRatioAero}
veces menor que esa reacción. En un motor real la corriente se satura y la
aceleración del rotor queda limitada por el par disponible; el modelo de
primer orden no tiene esa saturación, de modo que su $\dot\Omega$ inicial es
físicamente inalcanzable.

Aplicar $\dot{\vect h}_r$ con ese $\dot\Omega$ generaría golpes de guiñada
de casi 3~N·m (ocho veces el momento máximo de guiñada del vehículo,
$\val{taumaxY}$~N·m) que el dron real no siente. Modelarlo bien exige un
modelo eléctrico del motor con límite de corriente, que queda como trabajo
futuro (\cref{sec:limitaciones}). El modelo de Quan \cite{qq17} (cap.~6) y
su implementación CopterSim \cite{cs19}, de donde provienen los parámetros
del PAE-X450, tampoco incluyen este término: la omisión es coherente con la
fuente de los datos.

\subsection{Dinámica de los rotores}

Cada rotor sigue un primer orden hacia la velocidad de régimen que impone su
ESC (\cref{sec:propulsion}):
\begin{equation}
T_m\,\dot\Omega_i = \eta_i\,\omega_{ss}(\sigma_i) - \Omega_i ,
\label{eq:rotor}
\end{equation}
con $\eta_i\in(0,1]$ la eficiencia del motor (1 salvo en el escenario de falla,
\cref{sec:perturbaciones}).

\subsection{Fuerza específica y contacto con el suelo}

El acelerómetro ideal mide la fuerza específica, que el código calcula a
partir de la derivada del estado (\cod{specific\_force\_B}):
\begin{equation}
\vect f_{acc} = \R\T\big(\dot{\vect v} - g\,\eu\big)
= \frac{1}{m}\Big(\vect f_B + \R\T\vect F_{ext}\Big).
\label{eq:fesp}
\end{equation}
La segunda igualdad, que sale de \eqref{eq:newton}, tiene una consecuencia
que reaparece en la \cref{sec:estimadores}: el acelerómetro de un
multirrotor mide empuje y arrastre, nunca la gravedad. En hover inclinado y
sin velocidad mide $(0,0,-g)$ \emph{independientemente} de la inclinación.

El suelo se trata como una restricción unilateral simple
(\cod{Simulation.\_ground\_contact}): si $z>0$ se fija $z=0$ y se anula la
velocidad vertical descendente; si además el empuje total es menor que el
peso, el dron queda apoyado en sus patas: velocidad horizontal y angular
nulas y actitud nivelada conservando el rumbo. No hay fricción ni rebote:
basta para despegar y aterrizar, no para estudiar aterrizajes.
